\documentclass[11pt]{amsart}
\usepackage{calc,amssymb,amsmath,ulem,stmaryrd,fullpage}
\usepackage{enumerate}
\usepackage{alltt}
\RequirePackage[dvipsnames,usenames]{color}
\let\mffont=\sf
\normalem
%\input{kmacros3.sty}
\input{xy}
\xyoption{all}
\usepackage{hyperref}
\usepackage{graphicx}
\usepackage[all,cmtip]{xy}
\usepackage{verbatim}
\usepackage[left=0.95in,top=0.95in,right=0.95in,bottom=0.95in]{geometry}
\newcommand{\tauCohomology}{T}
\newcommand{\FCohomology}{S}


%\theoremstyle{theorem}
%\newtheorem*{mainthm*}{Main Theorem}

\newcommand{\note}[1]{\marginpar{\sffamily\tiny #1}}

\def\todo#1{\textcolor{Mahogany}%
{\sffamily\footnotesize\newline{\color{Mahogany}\fbox{\parbox{\textwidth-15pt}{#1}}}\newline}}

\renewcommand{\O}{\mathcal O}
\newcommand{\ie}{{\itshape i.e.} }
\newcommand{\roundup}[1]{\lceil #1 \rceil}
\newcommand{\rounddown}[1]{\lfloor #1 \rfloor}
\renewcommand{\to}[1][]{\xrightarrow{\ #1\ }}
\newcommand{\onto}[1][]{\protect{\xrightarrow{\ #1\ }\hspace{-0.8em}\rightarrow}}
\newcommand{\into}[1][]{\lhook \joinrel \xrightarrow{\ #1\ }}
%\newcommand{DBDual}[1]{{\underline{\omega}_{#1}}}
\newcommand{\DBDual}[1]{{{\uuline \omega}{}^{\mydot}_{#1}}}
\newcommand{\Tt}{{\mathfrak{T}}}
%\newcommand{\bn}{{\mathfrak{n}} }
\newcommand{\sol}[1]{ \vskip 6pt{\color{blue} {\bf Solution:} #1} \vskip 6pt}

\newcommand{\bR}{{\mathbb{R}}}
\newcommand{\va}{{\vec{a}}}
\newcommand{\vb}{{\vec{b}}}
\newcommand{\vzero}{{\vec{0}}}
\newcommand{\vi}{{\vec{i}}}
\newcommand{\vj}{{\vec{j}}}
\newcommand{\vk}{{\vec{k}}}
\newcommand{\vh}{{\vec{h}}}
\newcommand{\vr}{{\vec{r}}}
\newcommand{\vu}{{\vec{u}}}
\newcommand{\vv}{{\vec{v}}}
\newcommand{\vw}{{\vec{w}}}
\newcommand{\vx}{{\vec{x}}}
\newcommand{\vy}{{\vec{y}}}
\newcommand{\vz}{{\vec{z}}}
\newcommand{\vT}{{\vec{T}}}
\newcommand{\vN}{{\vec{N}}}
\newcommand{\vF}{{\vec{F}}}
\newcommand{\vG}{{\vec{G}}}
\newcommand{\bF}{\mathbb{F}}
\newcommand{\bQ}{\mathbb{Q}}
\newcommand{\bZ}{\mathbb{Z}}
\newcommand{\Inn}{\mathrm{Inn}}
\newcommand{\Aut}{\mathrm{Aut}}

\begin{document}
\title{HW \#3 -- Math 6310\\Fall 2019}
\author{Due: Friday, October 4th}
\maketitle
{\bf Definition 1.}  Suppose $R$ is a commutative integral domain.  A subset $W \subseteq R$ containing $1$ and not containing $0$ is called a \emph{multiplicative set} if for any $a,b \in W$, we have $ab \in W$.  (Note the conditions on 1 and 0 are not always assumed in the literature, we make them here for convenience).

In this case, we define $W^{-1}R$ to be the set of pairs $(a, b)$ such that $a \in R$ and $b \in W$, modulo the equivalence relation where $(a_1, b_1) \sim (a_2, b_2)$ if $a_1 b_2 = a_2 b_1$.  The equivalence class of $(a,b)$ we denote by $a/b$.  

$W^{-1}R$ forms a ring under the familiar rules for adding and multiplying fractions (you may take this as given).  

Furthermore, the map 
\[
    \phi : R \to W^{-1} R
    \]
sending $r \mapsto r/1$ is a ring homomorphism (you may also take it as given).

\begin{enumerate}
    \item Suppose that $R$ is a commutative integral domain, $W \subseteq R$ is a multiplicative set and $Q \subseteq R$ is a prime ideal such that $W \cap Q = \emptyset$.  Show that the set
    \[
    W^{-1} Q := \{ x/w \;|\; x \in Q, w \in W \} \subseteq W^{-1} R
    \]
    is a prime ideal of $W^{-1} R$.
    \vskip 3pt
    \emph{Hint:}  We checked this in class.
    \newpage
    \item  Suppose that $R$ is a
    commutative integral domain and $W \subseteq R$ is a multiplicative set.  Suppose that $P \subseteq W^{-1} R$ is a prime ideal and set $Q = P \cap R$ (where the intersection is taken with the isomorphic copy of $R \subseteq W^{-1}(R)$).  Show that $Q \cap W = \emptyset$ and that $W^{-1} Q = P$.
    \vskip 9pt
    Conclude that the prime ideals of $W^{-1} R$ are in bijection with the prime ideals $Q$ in $R$ such that $W \cap Q = \emptyset$.
    \newpage
    \item  Suppose that $W \subseteq R$ is a multiplicative set in a commutative integral domain, and $W$ is made up of only units (= invertible elements) of $R$.  Show that the map $\phi : R \to W^{-1} R$ from the first page is an isomorphism of rings.
    \newpage
    \item  Suppose that $R$ is a commutative integral domain, $W \subseteq R$ is a multiplicative set and $\psi : R \to S$ is a ring homomorphism between commutative rings.  Further suppose that for each $w \in W$, there is some $s \in S$ such that $s \psi(w) = 1_S$.  Show that one can factor $\psi$ as follows
    \[
    \xymatrix{
    R \ar[dr]_-{\phi} \ar[rr]^{\psi} & & S\\
    & W^{-1} R \ar@{.>}[ur]_-{\alpha}
    }
    \]
    and that the map $\alpha$ is unique.  This is called the \emph{universal property of localization}.
    \newpage
    \item Suppose that $R$ is a commutative integral domain and pick a nonzero $r \in R$.  Set $W = \{1, r, r^2, r^3, \dots\}$.  Show that we have an isomorphism between rings $R[x]/\langle rx - 1 \rangle \cong W^{-1}R$.  
    \newpage
    \item Show that any finite integral domain is a division ring.
    \newpage
    \item Do either of these distributive laws, for ideals $I, J, K$ in a commutative ring $R$ hold?
    \[
        \begin{array}{c}
            I\cdot (J + K) = I \cdot J + I \cdot K\\
            I \cap (J + K) = I \cap J + I \cap K
        \end{array}
    \]
    \newpage
    \item Suppose $R$ is a commutative ring and $I, J \subseteq R$ are ideals such that $I + J = R$ (in this case we say the ideals are (strongly) relatively prime).  Prove that $I \cap J = I \cdot J$.  What does this say for ideals in $\bZ$?
    \newpage
    \item Suppose that $R$ is a ring with ideals $I_1, I_2, \dots, I_k$ ideals such that $I_i + I_j = R$ for $i \neq j$.  Show that the system of equations 
    \[
        \begin{array}{rcl}
            x + I_1 & = & a_1 + I_1 \\
            x + I_2 & = & a_2 + I_2 \\
            \dots & = & \dots \\
            x+ I_k & = & a_k + I_k\\
        \end{array}
    \]
    always has a solution.  \vskip 3pt
    \emph{Hint:} First prove the case when $k = 2$ and show that two solutions differ by an element of $I_1 \cap I_2$.  You can then do induction and at some point you'll also want to show that $I_j + \bigcap_{i \neq j} I_i = R$.  
    \newpage
    \item Suppose that $I_1, I_2, \dots, I_k \subseteq R$ are ideals such that $I_i + I_j = R$ for $i \neq j$.  Prove that we have an isomorphism:
        \[
            R/(I_1 \cap I_2 \cap \dots \cap I_k) \cong (R/I_1) \oplus (R/I_2) \oplus \dots \oplus (R/I_k)
        \]        
\end{enumerate}

\end{document}

